\subsection{Data Preprocessing Pipeline}
\label{subsec:data_preprocessing}

The preprocessing stage uses one shared pipeline with dataset-specific choices
recorded in the \texttt{DatasetConfig} registry. The raw data are normalized,
the target is cleaned and encoded as binary, and each dataset is split into a
development partition and an untouched test partition before any learned
transformation is fitted. The final experiment settings are read from
\texttt{config.yaml}; dataset metadata remain in \texttt{src/config.py}.

\textbf{Dataset-specific cleaning.} Dataset 3 removes exact duplicate rows
because its mostly binary and ordinal representation produces repeated response
patterns. Dataset 2 converts zero values in \texttt{Cholesterol} and
\texttt{RestingBP} to missing values because zero is not physiologically
plausible for these measurements. Missing target values are removed before the
binary target is encoded.

\textbf{Learned transformations.} Numeric values are median-imputed and
standardized; categorical values are mode-imputed, binary-encoded, or one-hot
encoded. Dataset 2 additionally uses training-only IQR clipping on its
continuous clinical measurements. A variance filter and correlation filter are
then fit on processed training features. These transformations are fitted on
fold-training data inside cross-validation and on the full development split
for the final hold-out evaluation; the test rows never determine fitted
statistics.

Class imbalance is handled at model-fitting time with balanced sample weights,
not by resampling the data during preprocessing. This preserves the natural
class prevalence for validation and test metrics. The final processed shapes
and row counts are summarized in Table~\ref{tab:processed_dataset_summary}.

\begin{table}[H]
\centering
\begin{minipage}{0.90\linewidth}
\centering
	\refstepcounter{table}
	\begin{captionbelowboxaivn}
		\captable{Final processed train/test row counts and feature counts, after feature reduction, for the three datasets.}
		\label{tab:processed_dataset_summary}
	\end{captionbelowboxaivn}

	\renewcommand{\arraystretch}{1.1}
	\small

	\begin{tabularx}{\linewidth}{@{}
	>{\raggedright\arraybackslash}X
	>{\centering\arraybackslash}X
	>{\centering\arraybackslash}X
	>{\centering\arraybackslash}X@{}}
\textbf{Dataset} & \textbf{Train rows} & \textbf{Test rows} & \textbf{Final features (excl. target)} \\
\hline
Dataset 1 & 353{,}653 & 88{,}414 & 131 \\
Dataset 2 & 734 & 184 & 18 \\
Dataset 3 & 183{,}824 & 45{,}957 & 22 \\
	\end{tabularx}
\end{minipage}
\end{table}

The processed train/test files and the fold manifest are persisted under
\texttt{data/processed/<dataset>/}. The same fold manifest is reused by the
model and SHAP runners, while the fold-local fitting procedure is described in
the Model Experimental Design and SHAP Methodology subsections.
